
\begin{claim}{AC.35}{Proved partial reformulation; a finite etale torsor in the scalar case}
Let \(C/\mathbf F_q\) be a genus-two curve in characteristic
two, with a rational base point and the corresponding point-rigidified
Poincare family of degree-zero line bundles, defined over the base field. Suppose arithmetic Frobenius
on \(J_C\) is the scalar \([s]\), where \(s^2=q\).
On a nonempty rational open of \(J_C\), choose a rational basis
\(e_a\) of its one-dimensional nontrivial character cohomology,
and write the \(q\)-power coefficient as
\[
                     F_{\rm abs}^{\log_2q}e_a
                              =A(a)e_{[q]a}.
\]
For the models of AC.27, \(A\) is given explicitly by its
finite Cartier--Mumford product.
On the open \(U\) where this basis and \(A\) are invertible,
\[
                         \mathcal Y:\quad z^{q-1}=A
\]
is a finite etale \(\mu_{q-1}\)-torsor.
Its isomorphism class is independent of the basis on common
domains of definition.

If \(a\in U(\mathbf F_{q^m})\), the linear coherent return
scalar for the pointed character defined by \(a\) is
\[
                    \lambda(a)=
                    N_{\mathbf F_{q^m}/\mathbf F_q} A(a).
\]
This is exactly the arithmetic Frobenius value of
\(\mathcal Y\) at \(a\), with degree \(m\).
If \(a\) has odd prime order \(\ell\) and every nonzero multiple
of \(a\) lies in \(U\), its pointed cyclic cover over
\(\mathbf F_{q^m}\) satisfies (R) exactly when all these
Frobenius values differ from one.

The torsor furnished by \(C:y^2+y=x^5\) over \(\mathbf F_{16}\)
is geometrically nontrivial. This construction supplies a
finite etale object detecting the single-character return
condition, but does not prove that its nontrivial Frobenius
locus contains a whole nonzero prime cyclic subgroup.
\end{claim}
\begin{proof}
The scalar hypothesis gives the identity of Jacobian
endomorphisms
\[
                             [q]=F_J^2.
\]
Since the chosen rational functions are defined over
\(\mathbf F_q\), one has \(A(F_Ja)=A(a)^q\).
For a point \(a\) rational over \(\mathbf F_{q^m}\),
iterating the absolute Frobenius formula gives
\[
\begin{aligned}
 A_m(a)
  &=\prod_{j=0}^{m-1}
      A([q]^j a)^{q^{m-1-j}}\\
  &=A(a)^{\,q^{m-1}(1+q+\cdots+q^{m-1})}\\
  &=N_{\mathbf F_{q^m}/\mathbf F_q}A(a).
\end{aligned}
\]
The last expression lies in \(\mathbf F_q\), so the inverse
coefficient-field conjugation in AC.27 does not alter it.
Also \(F_J^m(a)=a\) implies \([q^m]a=a\).
If \(a\) has order \(n\), then \(n\mid s^m-1\), so
\(\mu_n\subset\mathbf F_{q^m}\).
The point-normalized line bundle therefore defines the
constant pointed cyclic cover over this same field.

For a different rational basis \(e'_a=b(a)e_a\), the coefficient
changes to
\[
 A'(a)=A(a)\frac{b(a)^q}{b([q]a)}
       =A(a)b(a)^{q-q^2}
       =A(a)\bigl(b(a)^{-q}\bigr)^{q-1}.
\]
This is precisely a change of Kummer generator, proving
basis independence. The relevant open is defined over
\(\mathbf F_q\) and is stable under \(F_J\), whose underlying
map is Frobenius; thus the return computation does not
require removing a new infinite collection of divisors.

Over a closed point with residue field \(\mathbf F_{q^m}\),
arithmetic Frobenius sends a root \(z\) to \(z^{q^m}\), so
its ratio is
\[
             z^{q^m-1}
              =A(a)^{(q^m-1)/(q-1)}
              =N_{\mathbf F_{q^m}/\mathbf F_q}A(a).
\]
This proves the exact comparison. For a prime-order
point, all nonzero multiples have the same field of
definition: multiplication by any \(1\le i<\ell\)
has an inverse on that cyclic subgroup.
Applying AC.17 or AC.27 to all its nontrivial characters
gives the stated criterion.

For the example of AC.31, \(q=16\), and the degree-three
closed point \(a\) has return value a nontrivial cube
root in \(\mathbf F_4^\times\).
The program now independently checks that each
\(4096\)-power scalar is the \(4096/16\) norm of its
\(16\)-power Cartier coefficient.
If the \(\mu_{15}\)-torsor were geometrically trivial,
its arithmetic Frobenius at a degree-three point would
be the cube of an element of \(\mu_{15}\), hence have
order dividing five. The observed order three is
impossible. Thus the torsor has nontrivial geometric
monodromy.

For general nonscalar Jacobian Frobenius the corresponding
periodicity condition is \(V_q^m(a)=a\), since on prime-to-two
torsion \(\sigma_{q^m}(a)=q^m a\) is equivalent to that
condition for the \(q\)-Verschiebung.
The simplification \([q]=F_J^2\), essential to the norm
and the Kummer gauge identity above, is then unavailable.
Even in the scalar case, ordinary Chebotarev applied to
\(U\) controls closed points of \(U\); it does not impose
prime order on \(a\) and simultaneous avoidance of one
on every nonzero multiple. Those are the precise additional
assertions still needed for A3.
\end{proof}

